\chapter{Limitations, Future Work and Conclusion}\label{ch:conclusion}

\section{Limitations}

\subsection{Open issues in the SDK and the HTTP mapping}

The socket-style SDK and the HTTP mapping were deliberately left out of the version~8 pass, which concentrated
on the protocol core. Their open findings are listed in Table~\ref{tab:open}. The first item in the table is
the most serious: an application that uses the SDK with encryption is still exposed to the GCM authentication-key
recovery described in Section~\ref{sec:nonce}, and fixing the SDK is the first item of future work.

\begin{table}[H]
\centering\small
\caption{Findings still open in the SDK and the HTTP mapping.}\label{tab:open}
\begin{tabularx}{\linewidth}{c>{\raggedright\arraybackslash}X}
\toprule
\textbf{\#} & \textbf{Issue} \\
\midrule
1 & SDK retransmissions re-seal with the retransmission bit in the AAD: the GCM nonce-reuse flaw remains in the SDK \\
2, 3 & SDK control messages are unauthenticated, and the SDK re-anchors its peer on any ACK, NACK or CLOSE \\
4 & The SDK server never authenticates the client, and allocates a connection, a socket and a thread per HELLO \\
12 & Idle SDK connections have no liveness check (HEARTBEAT is never sent) \\
13 & The SDK is stop-and-wait per message: only one message is in flight per direction \\
26 & About 17~MB per SDK connection: registry slots embed 16~KiB payloads while the SDK sends 1,200-byte blocks \\
27 & \code{fuse_close} frees the connection after two seconds even if a caller is still inside a call \\
28 & \code{fuse_connect} takes the PSK through \code{strlen}, so a key containing a zero byte is truncated \\
29 & HTTP: header injection (no CR/LF validation), timeout scope, reason truncation, detached threads, server status, and no tests \\
\bottomrule
\end{tabularx}
\end{table}

\subsection{Limitations of the protocol core}

\begin{itemize}
  \item \textbf{Fairness between lanes.} Every lane runs its own congestion controller, so a transfer with $N$
    lanes competes with other traffic on a shared bottleneck like $N$ TCP flows. Per-stream congestion control
    is the premise of the design, and it is what keeps a bulk stream from inflating the latency of a small one.
    The flip side is that a bulk transfer's \emph{aggregate} share should probably be coupled across its lanes.
  \item \textbf{Consecutive ports.} Bulk-transfer lanes need $N$ consecutive open UDP ports. A single-port
    transfer requires reworking the receive pipeline to demultiplex lanes by stream identifier on one socket.
  \item \textbf{Stragglers.} Static sharding means the slowest lane decides when the transfer completes; under
    heavy loss one lane can lag behind the others. Dynamic re-sharding or work stealing between lanes would
    address this.
  \item \textbf{Stage~4 flags are not exposed.} Unreliable, unordered and coalescing streams exist as tested
    primitives, and \code{fuse_bench} exercises a mixed session. No public transport offers them yet, so the
    mixed loss-tolerant and loss-intolerant session that motivates the protocol is not yet available to
    applications.
  \item \textbf{Forward secrecy and passwords.} The PSK path provides no forward secrecy. The DTLS~1.2 ECDHE-PSK
    component that would provide it is not wired into any transport. Accepting human-chosen passwords safely
    would need a PAKE such as CPace.
  \item \textbf{The multiplexed transfer} has no encryption, and therefore no control-plane authentication, and
    no address validation.
  \item \textbf{Platforms.} macOS uses one system call per datagram and cannot query the path MTU, so blocks do
    not grow there.
  \item \textbf{Evaluation.} All measurements are on one host with emulated impairments. There is no physical
    two-host run and no run on a real 1,500-MTU path, and the QUIC comparison predates version~8.
\end{itemize}

\section{Future work}

In order of priority, the next steps are:
\begin{enumerate}
  \item \textbf{Apply the version~8 hardening to the SDK}, reusing the bulk transfer's components rather than
    re-implementing them:
    \begin{itemize}
      \item the AAD fix;
      \item control-plane GMAC;
      \item path validation;
      \item Retry-style client validation before allocating per-connection state;
      \item heartbeats;
      \item a pipelined window instead of stop-and-wait;
      \item registry slots sized to the block.
    \end{itemize}
  \item \textbf{Fix the HTTP mapping's input validation}, and add tests.
  \item \textbf{Expose per-stream flags in a public API}, so an application can open a loss-tolerant and a
    loss-intolerant stream in one session. This is the feature that motivates the protocol.
  \item \textbf{Make the bulk transfer use a single port}, demultiplexing by stream identifier, and add work
    stealing between lanes.
  \item \textbf{Offer a forward-secret mode} by wiring the existing DTLS ECDHE-PSK handshake into session key
    establishment, keeping the per-lane AEAD data path. Add a PAKE for password-based deployments.
  \item \textbf{Couple congestion control across the lanes of one transfer}, so that the transfer as a whole is
    fair to other flows while its streams remain isolated from one another.
  \item \textbf{Evaluate on real networks}: two physical hosts, a 1,500-MTU wide-area path, and a fresh QUIC
    comparison against version~8.
  \item \textbf{Add fuzz testing of the decoders.} Every decoder is a natural fuzzing target, and none is fuzzed
    today.
\end{enumerate}

\section{Conclusion}

Fuse demonstrates that a transport whose reliability, ordering and congestion control are properties of
individual streams can be built on UDP with modest code. Its central ingredients are:
\begin{itemize}
  \item a ring-indexed selective-acknowledgement window;
  \item reorder-tolerant negative acknowledgements;
  \item per-stream AIMD with slow start, RFC~6298 timers and pacing;
  \item batched UDP I/O;
  \item per-lane parallel AEAD.
\end{itemize}
Earlier measurements showed the latency benefit of per-stream congestion control under mixed load, and
throughput competitive with a mature QUIC implementation.

The security and correctness review of version~7 showed how easily a datagram protocol goes wrong in ways that
tests of the happy path never reveal:
\begin{itemize}
  \item an innocuous flag in the AAD turned every retransmission into a key-recovery oracle;
  \item a convenient address re-anchoring became a hijacking primitive;
  \item an unused bitmask cost a factor of four in throughput under loss.
\end{itemize}
Protocol version~8 fixes all of these in the protocol core, mostly by adopting mechanisms whose design rationale
is well established in QUIC and TCP: Retry tokens, path validation, AEAD usage limits, selective
acknowledgement, RFC~6298 timers, pacing and MTU discovery. It does so without any loss of loopback throughput,
and with large gains under loss. The remaining work, above all the SDK, is now well defined.
